%%%%%%%%%%%%%%%%%%%%%%%%%%%%%%%%%%%%%%%%%%%%%%%%%%%%%%%%%%%%%%%%%%%%%%%%%%%%%%%
%  SÍLABO DE ECONOMÍA FINANCIERA — Plan de estudio integrado
%
%  Texto base:
%    SA  — Saunders, Cornett y Erhemjamts, Financial Markets and
%          Institutions (McGraw-Hill, 8.ª ed., 2022)
%  Valoración de activos:
%    CA  — Campbell, Financial Decisions and Markets (Princeton, 2017)
%    CO  — Cochrane, Asset Pricing (Princeton, ed. revisada, 2005)
%  Instituciones, banca e internacional:
%    ME  — Mishkin y Eakins, Financial Markets and Institutions
%          (Pearson, 10.ª ed. global, 2023)
%    GTB — Greenbaum, Thakor y Boot, Contemporary Financial
%          Intermediation (Academic Press, 4.ª ed., 2019)
%    KO  — Krugman, Obstfeld y Melitz, Economía internacional
%          (Pearson, 10.ª ed., 2016)
%  Programas PUCP integrados: 1ECO06 (Introducción a las Finanzas),
%    1FIN27 (Finanzas Internacionales), 1FIN31 (Regulación Financiera),
%    1FIN32 (Temas en Finanzas 1), 1FIN09 (Teoría y Gestión de Portafolios);
%    además, 1F0124 (Sistemas Financieros, Universidad del Pacífico).
%
%  Compilar con:  pdflatex EconomiaFinanciera.tex   (dos veces, por el índice)
%%%%%%%%%%%%%%%%%%%%%%%%%%%%%%%%%%%%%%%%%%%%%%%%%%%%%%%%%%%%%%%%%%%%%%%%%%%%%%%

\documentclass[11pt,a4paper]{article}

\usepackage[T1]{fontenc}
\usepackage[utf8]{inputenc}
\usepackage[spanish,es-noquoting,es-tabla]{babel}
\usepackage{lmodern}
\usepackage[scaled=0.92]{helvet}        % títulos en sans serif, al estilo McGraw-Hill
\usepackage[margin=2.4cm,headheight=15pt]{geometry}
\usepackage[table]{xcolor}
\usepackage{enumitem}
\usepackage{titlesec}
\usepackage{titletoc}
\usepackage{tabularx}
\usepackage{booktabs}
\usepackage{longtable}
\usepackage{amsmath}
\usepackage{amssymb}
\usepackage{fancyhdr}
\usepackage{tikz}
\usepackage[skins]{tcolorbox}

%------------------------------------------------------------------ colores ---
% Gama de azules
\definecolor{mhnavy}{RGB}{0,38,84}        % azul marino
\definecolor{mhdeep}{RGB}{0,59,119}       % azul profundo
\definecolor{mhblue}{RGB}{0,92,169}       % azul principal
\definecolor{mhroyal}{RGB}{35,85,190}     % azul real
\definecolor{mhcerulean}{RGB}{0,123,191}  % azul cerúleo
\definecolor{mhcyan}{RGB}{0,150,199}      % azul cian
\definecolor{mhsky}{RGB}{90,170,228}      % azul cielo
\definecolor{mhsteel}{RGB}{70,110,155}    % azul acero
\definecolor{mhslate}{RGB}{88,108,140}    % azul pizarra
\definecolor{mhgraphite}{RGB}{45,60,85}   % azul grafito (texto secundario)
\definecolor{mhice}{RGB}{210,229,247}     % azul hielo
\definecolor{mhpale}{RGB}{235,243,252}    % azul pálido
\definecolor{mhmist}{RGB}{246,250,255}    % azul bruma

\colorlet{azulosc}{mhnavy}
\colorlet{grisclaro}{mhpale}
% Un azul distinto para cada libro
\definecolor{colorSA}{RGB}{0,70,140}      % cobalto
\definecolor{colorCA}{RGB}{0,120,190}     % cerúleo
\definecolor{colorCO}{RGB}{0,135,165}     % azul petróleo
\definecolor{colorME}{RGB}{40,70,170}     % azul real
\definecolor{colorGTB}{RGB}{85,85,175}    % azul lavanda
\definecolor{colorKO}{RGB}{25,110,215}    % azul brillante
\definecolor{colorPUCP}{RGB}{85,105,135}  % azul pizarra

\usepackage[colorlinks=true,linkcolor=mhnavy,urlcolor=mhblue]{hyperref}

%------------------------------------------------------------------ estilos ---
\titleformat{\section}[block]
  {\sffamily\LARGE\bfseries\color{mhnavy}}{\thesection.}{0.6em}
  {\textcolor{mhblue}{\rule[-0.1ex]{0.55em}{0.85em}}\hspace{0.45em}}
  [\vspace{-0.4ex}{\color{mhblue}\titlerule[1.6pt]}\vspace{1.2pt}{\color{mhsky}\titlerule[0.6pt]}]
\titlespacing*{\section}{0pt}{2.6ex plus 1ex minus .2ex}{1.6ex plus .2ex}

\titleformat{\subsection}
  {\sffamily\large\bfseries\color{mhblue}}{\thesubsection}{0.6em}
  {\textcolor{mhsky}{\rule[0.1ex]{0.45em}{0.45em}}\hspace{0.45em}}
\titlespacing*{\subsection}{0pt}{2.2ex plus .8ex minus .2ex}{1ex plus .2ex}

% Índice
\titlecontents{section}[0pt]
  {\addvspace{3pt}\sffamily\bfseries\color{mhnavy}}
  {}{\textcolor{mhcerulean}{\rule[0.15ex]{0.4em}{0.4em}}\hspace{0.6em}}
  {{\color{mhsky}\titlerule*[0.7em]{.}}\textcolor{mhblue}{\contentspage}}

% Cabecera y pie
\pagestyle{fancy}
\fancyhf{}
\lhead{\small\sffamily\bfseries\color{mhnavy}Sílabo de Economía Financiera}
\rhead{\footnotesize\sffamily\color{mhblue}Saunders et al. $\cdot$ Campbell $\cdot$ Cochrane $\cdot$ Mishkin--Eakins $\cdot$ Greenbaum et al.}
\cfoot{{\setlength{\fboxsep}{3pt}\colorbox{mhblue}{\makebox[1.8em]{\small\sffamily\bfseries\color{white}\thepage}}}}
\renewcommand{\headrule}{%
  {\color{mhblue}\hrule height 1.1pt width\headwidth}\vskip 1.2pt
  {\color{mhsky}\hrule height 0.4pt width\headwidth}\vskip -2.7pt}

% Listas
\setlist[itemize,1]{label=\textcolor{mhblue}{\scriptsize$\blacksquare$}}
\setlist[enumerate,1]{label=\textcolor{mhblue}{\sffamily\bfseries\arabic*.}}

% Tablas: cabecera azul marino y filas alternas en azul pálido
\arrayrulecolor{mhblue}
\setlength{\aboverulesep}{0pt}
\setlength{\belowrulesep}{0pt}
\renewcommand{\arraystretch}{1.3}
\newcommand{\cab}[1]{\textcolor{white}{\sffamily\bfseries #1}}
\newcommand{\filasalternas}{\rowcolors{2}{white}{mhpale}}

% Encabezados de párrafo en azul
\newcommand{\destacado}[1]{{\sffamily\bfseries\color{mhblue}#1}}

% Franja decorativa de azules
\newcommand{\franjas}[1][3pt]{%
  \noindent\begin{tikzpicture}
    \foreach \c [count=\i from 0] in {mhnavy,mhdeep,mhblue,mhroyal,mhcerulean,mhcyan,mhsky,mhice}
      \fill[\c] (\i*\linewidth/8,0) rectangle ++(\linewidth/8,#1);
  \end{tikzpicture}\par}

%---------------------------------------------------------------- comandos ----
% Etiqueta de libro: recuadro de color con la sigla
\newcommand{\etiqueta}[2]{%
  {\setlength{\fboxsep}{1.4pt}\colorbox{#1}{\sffamily\bfseries\color{white}\strut#2}}}
% Referencias abreviadas a cada libro
\newcommand{\SA}[1]{\etiqueta{colorSA}{SA}\nobreak\hspace{0.25em}\textcolor{colorSA}{#1}}
\newcommand{\CA}[1]{\etiqueta{colorCA}{CA}\nobreak\hspace{0.25em}\textcolor{colorCA}{#1}}
\newcommand{\CO}[1]{\etiqueta{colorCO}{CO}\nobreak\hspace{0.25em}\textcolor{colorCO}{#1}}
\newcommand{\ME}[1]{\etiqueta{colorME}{ME}\nobreak\hspace{0.25em}\textcolor{colorME}{#1}}
\newcommand{\GTB}[1]{\etiqueta{colorGTB}{GTB}\nobreak\hspace{0.25em}\textcolor{colorGTB}{#1}}
\newcommand{\KO}[1]{\etiqueta{colorKO}{KO}\nobreak\hspace{0.25em}\textcolor{colorKO}{#1}}

% Un tema: número en recuadro azul + título
\newcommand{\tema}[1]{\temaaux#1\relax}
\def\temaaux#1 #2\relax{%
  \par\addvspace{1.1em}\noindent
  {\sffamily\bfseries\large
   \hangindent=2.9em\hangafter=1
   \makebox[2.9em][l]{{\setlength{\fboxsep}{2.5pt}\colorbox{mhblue}{\color{white}#1}}}%
   {\color{mhnavy}#2}\par}%
  \nobreak\vspace{0.25em}\noindent{\color{mhsky}\rule{\linewidth}{0.6pt}}\par\nobreak\vspace{0.35em}}

% Línea de lecturas con marcador y etiqueta de color
\newcommand{\linealectura}[3]{%
  \par\nobreak\noindent\hspace{1.2em}%
  \begin{minipage}{0.92\textwidth}\small
    \textcolor{#1}{\rule[0.1ex]{0.45em}{0.45em}}\hspace{0.45em}%
    {\sffamily\bfseries\color{#1}#2}\enspace#3%
  \end{minipage}\par\vspace{0.3em}}
\newcommand{\lecturas}[1]{\linealectura{mhblue}{Lecturas:}{#1}}
\newcommand{\apoyo}[1]{\linealectura{mhroyal}{Profundización:}{#1}}
\newtcolorbox{cajanota}{enhanced, sharp corners, boxrule=0pt, frame hidden,
  colback=mhpale, borderline west={2.5pt}{0pt}{mhsky},
  left=8pt, right=6pt, top=4pt, bottom=4pt, before skip=0pt, after skip=0pt,
  fontupper=\small\color{mhgraphite}}
\newcommand{\nota}[1]{%
  \par\nobreak\noindent\hspace{1.2em}%
  \begin{minipage}{0.92\textwidth}\begin{cajanota}#1\end{cajanota}\end{minipage}\par\vspace{0.3em}}
\newcommand{\consultalib}[1]{\linealectura{mhslate}{Consulta:}{{\color{mhgraphite}#1}}}
% Artículos y textos tomados de los programas PUCP
\newcommand{\articulos}[1]{\linealectura{colorPUCP}{Artículos (PUCP):}{#1}}
% Artículos que no proceden de los programas (entran en las horas)
\newcommand{\articulosB}[1]{\linealectura{mhcerulean}{Artículos:}{#1}}
\newcommand{\consulta}[1]{\linealectura{colorPUCP}{Consulta (PUCP):}{{\color{mhgraphite}#1}}}

% Encabezado de bloque: banda azul marino, franja de azules y recuadro de presentación
\newcommand{\bloque}[2]{%
  \clearpage
  \phantomsection\addcontentsline{toc}{section}{#1}%
  \begin{tcolorbox}[enhanced, sharp corners, boxrule=0pt, frame hidden,
      colback=mhnavy, borderline west={7pt}{0pt}{mhsky},
      left=14pt, right=10pt, top=10pt, bottom=10pt, before skip=0pt, after skip=0pt]
    {\sffamily\Large\bfseries\color{white}#1}
  \end{tcolorbox}%
  \nointerlineskip\franjas[4pt]
  \vspace{0.8em}
  \begin{tcolorbox}[enhanced, sharp corners, boxrule=0pt, frame hidden,
      colback=mhpale, borderline west={3pt}{0pt}{mhblue},
      borderline south={0.6pt}{0pt}{mhsky},
      left=10pt, right=8pt, top=6pt, bottom=6pt, before skip=0pt, after skip=0pt,
      fontupper=\small\color{mhgraphite}]
    #2
  \end{tcolorbox}%
  \par\medskip}

%================================================================== documento ==
\begin{document}

\begin{titlepage}
\begin{tikzpicture}[remember picture,overlay]
  % banda superior
  \fill[mhnavy] (current page.north west) rectangle ([yshift=-5.4cm]current page.north east);
  \fill[mhsky] (current page.north west) rectangle ([xshift=1.1cm,yshift=-5.4cm]current page.north west);
  % franja de azules bajo la banda
  \foreach \c [count=\i from 0] in {mhnavy,mhdeep,mhblue,mhroyal,mhcerulean,mhcyan,mhsky,mhice}
    \fill[\c] ([xshift=\i*\paperwidth/8,yshift=-5.4cm]current page.north west)
      rectangle ++(\paperwidth/8,-0.35cm);
  % franja inferior
  \foreach \c [count=\i from 0] in {mhice,mhsky,mhcyan,mhcerulean,mhroyal,mhblue,mhdeep,mhnavy}
    \fill[\c] ([xshift=\i*\paperwidth/8]current page.south west)
      rectangle ++(\paperwidth/8,0.9cm);
\end{tikzpicture}%
\centering
\vspace*{0.5cm}
{\sffamily\Huge\bfseries\color{white} Sílabo de Economía Financiera\par}
\vspace{0.4cm}
{\sffamily\Large\color{mhice} Plan de estudio integrado\par}
\vspace{1.9cm}
{\sffamily\large\bfseries\color{mhblue}\textcolor{mhsky}{\rule[0.5ex]{2.2cm}{0.6pt}}\quad Texto base\quad\textcolor{mhsky}{\rule[0.5ex]{2.2cm}{0.6pt}}\par}
\vspace{0.3cm}
{\large\etiqueta{colorSA}{SA} \quad {\color{mhnavy}Anthony Saunders, Marcia Millon Cornett y Otgontsetseg Erhemjamts}\par}
{\normalsize \emph{\color{colorSA}Financial Markets and Institutions}\par}
{\small\color{mhslate} McGraw-Hill, 8.ª edición, 2022 --- 25 capítulos\par}
\vspace{0.6cm}
{\sffamily\large\bfseries\color{mhblue}\textcolor{mhsky}{\rule[0.5ex]{2.2cm}{0.6pt}}\quad Valoración de activos\quad\textcolor{mhsky}{\rule[0.5ex]{2.2cm}{0.6pt}}\par}
\vspace{0.3cm}
{\large\etiqueta{colorCA}{CA} \quad {\color{mhnavy}John Y. Campbell}\par}
{\normalsize \emph{\color{colorCA}Financial Decisions and Markets: A Course in Asset Pricing}\par}
{\small\color{mhslate} Princeton University Press, 2017 --- 12 capítulos\par}
\vspace{0.2cm}
{\large\etiqueta{colorCO}{CO} \quad {\color{mhnavy}John H. Cochrane}\par}
{\normalsize \emph{\color{colorCO}Asset Pricing}\par}
{\small\color{mhslate} Princeton University Press, edición revisada, 2005 --- 21 capítulos\par}
\vspace{0.6cm}
{\sffamily\large\bfseries\color{mhblue}\textcolor{mhsky}{\rule[0.5ex]{1.2cm}{0.6pt}}\quad Instituciones, banca y finanzas internacionales\quad\textcolor{mhsky}{\rule[0.5ex]{1.2cm}{0.6pt}}\par}
\vspace{0.3cm}
{\large\etiqueta{colorME}{ME} \quad {\color{mhnavy}Frederic S. Mishkin y Stanley G. Eakins}\par}
{\normalsize \emph{\color{colorME}Financial Markets and Institutions}\par}
{\small\color{mhslate} Pearson, 10.ª edición global, 2023 --- 27 capítulos\par}
\vspace{0.2cm}
{\large\etiqueta{colorGTB}{GTB} \quad {\color{mhnavy}Stuart I. Greenbaum, Anjan V. Thakor y Arnoud W. A. Boot}\par}
{\normalsize \emph{\color{colorGTB}Contemporary Financial Intermediation}\par}
{\small\color{mhslate} Academic Press, 4.ª edición, 2019 --- 18 capítulos\par}
\vspace{0.2cm}
{\large\etiqueta{colorKO}{KO} \quad {\color{mhnavy}Paul R. Krugman, Maurice Obstfeld y Marc J. Melitz}\par}
{\normalsize \emph{\color{colorKO}Economía internacional: teoría y política}\par}
{\small\color{mhslate} Pearson, 10.ª edición, 2016 --- capítulos de finanzas internacionales\par}
\vspace{0.2cm}
{\large\color{mhnavy} Xavier Freixas y Jean-Charles Rochet\par}
{\normalsize \emph{\color{mhsteel}Microeconomics of Banking}\par}
{\small\color{mhslate} MIT Press, 1999 --- capítulos de teoría bancaria\par}
\vspace{0.5cm}
{\normalsize\color{colorPUCP} Con temas y lecturas de cinco programas de la PUCP\par}
\vfill
{\setlength{\fboxsep}{6pt}\colorbox{mhblue}{\normalsize\sffamily\bfseries\color{white} Ocho bloques temáticos --- unos 5 meses}\par}
\vspace{0.2cm}
\end{titlepage}

\tableofcontents
\newpage

%------------------------------------------------------------------------------
\section*{Cómo usar este sílabo}
\addcontentsline{toc}{section}{Cómo usar este sílabo}

El plan está organizado en \textbf{ocho bloques temáticos} y avanza de lo
institucional a lo analítico y de vuelta a la política: primero el sistema
financiero, las tasas de interés y los mercados con Saunders et al.; después
la valoración de activos con Campbell y Cochrane, del portafolio de media y
varianza al factor de descuento estocástico y las macrofinanzas; luego la
banca, la regulación y las crisis; y, al final, las finanzas internacionales
con énfasis en el Perú. Conviene cerrar cada bloque antes de abrir el
siguiente.

\medskip
\noindent\destacado{Por qué un curso aparte.} El sílabo de macroeconomía deja
fuera las finanzas, la banca y casi toda la macroeconomía internacional; de la
política peruana solo trata la dolarización desde la política monetaria (su
tema 13.7). Este plan llena ese vacío de forma acotada, con el foco puesto en
la economía y no en la gestión de empresas o de inversiones.

\medskip
\noindent\destacado{Requisitos.} Del sílabo de microeconomía, la decisión bajo
incertidumbre (Bloque 4) y el equilibrio general con mercados contingentes
(Bloque 10); del de macroeconomía, el consumo intertemporal (Bloque 2), el
modelo de Ramsey (Bloque 4) y el dinero (Bloque 8); del de econometría, el GMM
(Bloque 9). El Bloque 1 puede empezarse sin ellos.

\medskip
\noindent\destacado{Criterio general de reparto} (en el orden en que se lee cada tema):
\begin{itemize}[leftmargin=1.4em,itemsep=2pt]
  \item \textcolor{colorSA}{\textbf{Saunders et al.}} aporta la mecánica, con
        ejemplos numéricos: valoración, duración, mercados e instrumentos,
        balances bancarios y gestión de riesgos.
  \item \textcolor{colorME}{\textbf{Mishkin y Eakins}} aporta la economía: de
        qué dependen las tasas de interés, por qué existen los intermediarios,
        por qué ocurren las crisis y cómo funciona el sistema financiero
        internacional. Es el texto de Sistemas Financieros (Universidad del
        Pacífico). En los Bloques 1, 2 y 6 cada tema se lee en uno de los dos,
        según sea de mecánica o de economía, y el otro queda como consulta.
  \item \textcolor{colorCA}{\textbf{Campbell}} es la lectura principal de
        valoración de activos (Bloques 3 a 5): riguroso, breve y conectado con
        la evidencia empírica.
  \item \textcolor{colorCO}{\textbf{Cochrane}} es la referencia del enfoque del
        factor de descuento estocástico y de su estimación por GMM. Se lee
        \emph{después} de Campbell.
  \item \textbf{Freixas y Rochet} dan la teoría microeconómica de la banca:
        por qué existen los intermediarios, las corridas y la regulación
        (Bloques 6 y 7). Se usa la edición en inglés de la carpeta de
        macroeconomía, un escaneo completo sin capa de texto.
  \item \textcolor{colorGTB}{\textbf{Greenbaum, Thakor y Boot}} cubren la
        titulización, el contrato de depósito y la crisis de 2008 (temas 6.4 y
        7.2); el resto queda como consulta.
  \item \textcolor{colorKO}{\textbf{Krugman, Obstfeld y Melitz}}, junto con
        las lecturas de 1FIN27, sostienen el Bloque 8.
\end{itemize}

\medskip
\noindent\destacado{Qué entra en las horas.} Las líneas \emph{Lecturas},
\emph{Profundización}, \emph{Artículos} y \emph{Artículos (PUCP)} entran en el
cómputo de horas; \emph{Artículos} son lecturas que no proceden de los
programas.
Las líneas \emph{Consulta} y \emph{Consulta (PUCP)} remiten a libros de la
carpeta y a referencias de los programas que no entran (Anexo E).

\medskip
\noindent\destacado{Material excluido.} Quedan fuera las finanzas corporativas,
la gestión de portafolios, la matemática financiera en tiempo continuo, la
valoración de derivados y la gestión de riesgos con fines profesionales. El
detalle está en el Anexo B.

\medskip
\noindent\destacado{Notación.} \SA{n} remite al capítulo $n$ de Saunders et al.,
y \CA{n}, \CO{n}, \ME{n}, \GTB{n} y \KO{n}, al capítulo de los demás textos.
Si solo entra parte de un capítulo se indican secciones: \CA{6.4} es la
sección 4 del capítulo 6 de Campbell. Saunders et al., Mishkin y Eakins y
Greenbaum et al.\ no numeran sus secciones, así que se citan por capítulo o
apéndice (\SA{3A}). \emph{Consulta} remite a libros de la carpeta que no
entran en el cómputo de horas (Anexo E).

%------------------------------------------------------------------------------
\newpage
\section*{Cronograma por semanas}
\addcontentsline{toc}{section}{Cronograma por semanas}

Calendario del plan: 200 horas distribuidas en \textbf{20 semanas de
contenido}, que equivalen a \textbf{unos 5 meses de calendario}, a un ritmo de
diez horas semanales.

\medskip
\noindent Todas las semanas son de materia nueva y no hay semanas de colchón.
La columna de horas es la que gobierna el reparto. Si hay semanas perdidas por
feriados, viajes o imprevistos, el ajuste correcto es alargar el calendario,
no comprimir los bloques.

\medskip
{\filasalternas
\noindent\begin{tabularx}{\textwidth}{lXcc}
\toprule
\rowcolor{mhnavy}\cab{Semanas} & \cab{Contenido} & \cab{Sem.} & \cab{Horas} \\
\midrule
1--2 & Bloque 1 --- El sistema financiero y las tasas de interés & 2 & 20 \\
3--4 & Bloque 2 --- Mercados e instrumentos financieros & 2 & 20 \\
5--6 & Bloque 3 --- Riesgo, portafolio y CAPM & 2 & 20 \\
\rowcolor{mhice}\textbf{7--9} & \textbf{Bloque 4 --- El factor de descuento estocástico} & \textbf{3} & \textbf{30} \\
\rowcolor{mhice}\textbf{10--12} & \textbf{Bloque 5 --- Valoración intertemporal y macrofinanzas} & \textbf{3} & \textbf{30} \\
\rowcolor{mhice}\textbf{13--15} & \textbf{Bloque 6 --- Intermediación financiera y banca} & \textbf{3} & \textbf{30} \\
\rowcolor{mhice}\textbf{16--18} & \textbf{Bloque 7 --- Regulación, crisis y estabilidad financiera} & \textbf{3} & \textbf{30} \\
19--20 & Bloque 8 --- Finanzas internacionales & 2 & 20 \\
\midrule
\rowcolor{mhice}\multicolumn{2}{l}{\textbf{Total}} & \textbf{20} & \textbf{200} \\
\bottomrule
\end{tabularx}
}

\medskip
\noindent\destacado{Ritmo supuesto.} Diez horas por semana, todas de contenido
nuevo: por ejemplo, cuatro sesiones de dos horas entre semana y una sesión de
dos horas el sábado para ejercicios y datos.

\medskip
\noindent\destacado{El tramo largo.} Los Bloques 4 y 5 (semanas 7--12) son el
núcleo analítico del plan y el punto donde el avance suele atascarse: el
paso del CAPM al factor de descuento estocástico exige el equilibrio con
mercados contingentes del sílabo de microeconomía.

\medskip
\noindent\destacado{Hitos de control.} Conviene verificar el avance al terminar las
semanas 6, 12 y 18. Si en alguno de esos cortes se acumula más de una semana de
retraso, se alarga el calendario y no se comprime el bloque siguiente.

\subsection*{Detalle semana a semana}

{\small\filasalternas
\noindent\begin{tabularx}{\textwidth}{clXc}
\toprule
\rowcolor{mhnavy}\cab{Sem.} & \cab{Bloque} & \cab{Temas} & \cab{Horas} \\
\midrule
1  & 1 & 1.1 Mercados e instituciones financieras; 1.2 Determinantes de las tasas de interés & 10 \\
2  & 1 & 1.3 Tasas de interés y valoración: duración y convexidad; 1.4 Banco central, política monetaria y tasas & 10 \\
\midrule
3  & 2 & 2.1 Mercados de dinero y de bonos; 2.2 Mercado de acciones y eficiencia & 10 \\
4  & 2 & 2.3 Hipotecas y titulización; 2.4 Derivados & 10 \\
\midrule
5  & 3 & 3.1 Decisión bajo incertidumbre (repaso); 3.2 Elección de portafolio de media y varianza & 10 \\
6  & 3 & 3.3 CAPM, APT y modelos de factores & 10 \\
\midrule
7  & 4 & 4.1 La ecuación básica de valoración & 10 \\
8  & 4 & 4.2 Mercados contingentes y ausencia de arbitraje; 4.3 Factor de descuento, betas y frontera & 10 \\
9  & 4 & 4.4 GMM y pruebas de modelos de valoración & 10 \\
\midrule
10 & 5 & 5.1 Valor presente, eficiencia y predictibilidad; 5.2 Valoración basada en el consumo y sus enigmas & 10 \\
11 & 5 & 5.3 Preferencias recursivas, riesgo de largo plazo y hábitos; 5.4 Valoración basada en la producción & 10 \\
12 & 5 & 5.5 Estructura temporal de tasas; 5.6 Aplicación: tasas, precios de activos e IA & 10 \\
\midrule
13 & 6 & 6.1 Por qué existen los intermediarios; 6.2 Bancos e instituciones financieras & 10 \\
14 & 6 & 6.3 Riesgos de las instituciones financieras; 6.4 Depósitos, corridas, titulización y banca en la sombra & 10 \\
15 & 6 & 6.5 El sistema financiero peruano & 10 \\
\midrule
16 & 7 & 7.1 Regulación bancaria: de Basilea I a Basilea III & 10 \\
17 & 7 & 7.2 Crisis financieras: anatomía e historia & 10 \\
18 & 7 & 7.3 Prestamista de última instancia y política macroprudencial & 10 \\
\midrule
19 & 8 & 8.1 Tipo de cambio y paridades; 8.2 El sistema monetario internacional & 10 \\
20 & 8 & 8.3 Crisis en economías emergentes; 8.4 Política monetaria y cambiaria en el Perú & 10 \\
\bottomrule
\end{tabularx}
}

%------------------------------------------------------------------------------
\newpage
\bloque{Bloque 1 --- El sistema financiero y las tasas de interés --- 2 semanas}
{Los conceptos de base de todo el plan: qué hacen los mercados y las
instituciones financieras, de qué dependen las tasas de interés y cómo se
valora un activo. Cubre las primeras semanas de 1ECO06 y de 1FIN32 y la
Unidad 1 de Sistemas Financieros (Universidad del Pacífico).}

\tema{1.1 Mercados e instituciones financieras}
\lecturas{\ME{1--2}}
\consultalib{\SA{1} \quad Cecchetti y Schoenholtz (2015), caps.~1--3}
\nota{Mercados primarios y secundarios, de dinero y de capitales; el papel de
los intermediarios y la globalización financiera. El apéndice de \SA{1} sobre
la crisis de 2008, de consulta, adelanta el Bloque 7.}

\tema{1.2 Determinantes de las tasas de interés}
\lecturas{\ME{4--5}}
\consultalib{\SA{2}}
\nota{La teoría de fondos prestables, la prima por riesgo de impago, liquidez
e inflación, y una primera mirada a la estructura temporal (que se formaliza
en el tema 5.5).}

\tema{1.3 Tasas de interés y valoración de activos: duración y convexidad}
\lecturas{\SA{3} \quad \SA{3A--3B}}
\consultalib{\ME{3} \quad Bodie, Kane y Marcus (2018), cap.~14}
\nota{Valoración de bonos y acciones por flujos descontados, efecto de la
madurez y del cupón, duración e inmunización.}

\tema{1.4 Banco central, política monetaria y tasas de interés}
\lecturas{\ME{9--10}}
\consultalib{\SA{4} \quad Cecchetti y Schoenholtz (2015), caps.~15--18}
\nota{Instrumentos de política monetaria y su efecto sobre las tasas de
mercado. Enlaza con el Bloque 8 del sílabo de macroeconomía; la aplicación
al Perú está en el tema 8.4.}

%------------------------------------------------------------------------------
\bloque{Bloque 2 --- Mercados e instrumentos financieros --- 2 semanas}
{Un recorrido por los mercados con Saunders et al.: qué se negocia, quién
participa y cómo se valora cada instrumento. Cubre 1FIN32 y la Unidad 3 de
Sistemas Financieros.}

\tema{2.1 Mercados de dinero y de bonos}
\lecturas{\SA{5--6}}
\consultalib{\ME{11--12} \quad Tuckman (2002), caps.~1--3}

\tema{2.2 Mercado de acciones y eficiencia de los mercados}
\lecturas{\SA{8} \quad \ME{6}}
\consultalib{\ME{13} \quad Bodie, Kane y Marcus (2018), caps.~11--12}
\nota{Hipótesis de mercados eficientes y sus anomalías; \SA{8B} explica los
estudios de eventos.}

\tema{2.3 Mercado hipotecario y titulización}
\lecturas{\SA{7} \quad \SA{25}}
\consultalib{\ME{14}}
\nota{Los valores respaldados por hipotecas, protagonistas de la crisis de
2008 (tema 7.2).}

\tema{2.4 Derivados: futuros, opciones y \emph{swaps}}
\lecturas{\SA{10} \quad \SA{24}}
\consultalib{Hull (2021), caps.~1--2 \quad Bodie, Kane y Marcus (2018), caps.~20--23}
\nota{Qué son y para qué los usan las instituciones financieras: cobertura de
riesgos. La valoración de opciones en tiempo continuo queda fuera
(\SA{10A} la resume).}

%------------------------------------------------------------------------------
\bloque{Bloque 3 --- Riesgo, portafolio y CAPM --- 2 semanas}
{La teoría clásica de 1ECO06 y de 1FIN09 en la versión de Campbell. Supone el
Bloque 4 del sílabo de microeconomía.}

\tema{3.1 Decisión bajo incertidumbre y aversión al riesgo (repaso)}
\consultalib{\CA{1}}
\nota{Utilidad esperada, aversión absoluta y relativa al riesgo y dominancia
estocástica, en el lenguaje de las finanzas. Repite el Bloque 4 del sílabo de
microeconomía, que es requisito; por eso no suma horas.}

\tema{3.2 Elección de portafolio de media y varianza}
\lecturas{\CA{2}}
\consultalib{Bodie, Kane y Marcus (2018), caps.~5--7}
\nota{Diversificación, frontera eficiente, activo libre de riesgo y teorema de
separación (Markowitz).}

\tema{3.3 CAPM, APT y modelos de factores}
\lecturas{\CA{3} \quad \SA{8A}}
\apoyo{\CO{9.1--9.4}}
\consultalib{Bodie, Kane y Marcus (2018), caps.~8--10 y 13}
\nota{El CAPM y su evidencia, los factores de Fama--French y la teoría de
valoración por arbitraje.}

%------------------------------------------------------------------------------
\bloque{Bloque 4 --- El factor de descuento estocástico --- 3 semanas}
{El enfoque moderno de la valoración de activos: un solo principio ---precio
igual a la esperanza del pago descontado por el factor de descuento
estocástico--- del que se derivan el CAPM, la ausencia de arbitraje y los
modelos de factores. Supone el Bloque 10 del sílabo de microeconomía.}

\tema{4.1 La ecuación básica de valoración}
\lecturas{\CO{1--2}}
\nota{El factor de descuento como tasa marginal de sustitución, la tasa libre
de riesgo, la corrección por riesgo y el equilibrio general.}

\tema{4.2 Mercados contingentes y ausencia de arbitraje}
\lecturas{\CA{4.1--4.2} \quad \CO{3--4}}
\nota{Probabilidades neutrales al riesgo, ley de un solo precio y existencia de
un factor de descuento positivo.}

\tema{4.3 Factor de descuento, betas y frontera de media y varianza}
\lecturas{\CA{4.3} \quad \CO{5.1--5.2} \quad \CO{6.1--6.3}}
\nota{La cota de Hansen--Jagannathan y la equivalencia entre factor de
descuento, representación beta y frontera eficiente.}

\tema{4.4 GMM y pruebas de modelos de valoración}
\lecturas{\CA{4.4--4.5} \quad \CO{10} \quad \CO{12}}
\articulosB{Hansen y Singleton (1982)}
\consultalib{Cochrane (2005), caps.~11 y 13 \quad Campbell, Lo y MacKinlay (1997)}
\nota{Estimación de ecuaciones de Euler por GMM, que inauguran Hansen y
Singleton (1982), y regresiones de Fama--MacBeth.
Aplica el Bloque 9 del sílabo de econometría. \CA{4.5} trata los límites al
arbitraje.}

%------------------------------------------------------------------------------
\bloque{Bloque 5 --- Valoración intertemporal y macrofinanzas --- 3 semanas}
{El puente entre la valoración de activos y la macroeconomía. Explica por qué
los precios de los activos contienen información sobre las expectativas de
crecimiento, y cierra con la aplicación a la inteligencia artificial.}

\tema{5.1 Valor presente, eficiencia y predictibilidad de los retornos}
\lecturas{\CA{5}}
\apoyo{\CO{20.1}}
\nota{La descomposición de Campbell--Shiller y la volatilidad de los precios
de las acciones.}

\tema{5.2 Valoración basada en el consumo y sus enigmas}
\lecturas{\CA{6.1--6.3}}
\apoyo{\CO{21.1}}
\articulosB{Lucas (1978) \quad Mehra y Prescott (1985)}
\nota{El enigma de la prima de riesgo, el de la tasa libre de riesgo y el de la
volatilidad. Usa la ecuación de Euler del Bloque 2 del sílabo de
macroeconomía. Lucas (1978) es el modelo de equilibrio del que parte la
valoración basada en el consumo, y Mehra y Prescott (1985), el artículo que
planteó el enigma de la prima de riesgo.}

\tema{5.3 Preferencias recursivas, riesgo de largo plazo y hábitos}
\lecturas{\CA{6.4--6.8}}
\apoyo{\CO{21.2}}
\nota{Epstein--Zin, riesgo de largo plazo, aversión a la ambigüedad y
formación de hábitos.}

\tema{5.4 Valoración basada en la producción}
\lecturas{\CA{7}}
\nota{Inversión con costos de ajuste, la $q$ de Tobin y el equilibrio general
con producción. Enlaza con el modelo de Ramsey del sílabo de macroeconomía.}

\tema{5.5 La estructura temporal de las tasas de interés}
\lecturas{\CA{8.1--8.4} \quad \CO{19.1--19.3}}
\consultalib{Bodie, Kane y Marcus (2018), cap.~15}
\nota{Hipótesis de expectativas, primas por plazo y modelos afines.}

\tema{5.6 Aplicación: tasas de interés, precios de activos e inteligencia artificial}
\lecturas{Chow, Halperin y Mazlish (2023) \quad Van Nieuwerburgh (2026)}
\nota{Más crecimiento esperado eleva las tasas reales (ecuación de Euler del
tema 5.2); el primer trabajo usa esa relación para leer qué esperan los
mercados de la IA. El segundo estima cuánto cuesta financiar la inversión en
IA y qué rentabilidad necesita. Ambos están en el módulo G.6 del sílabo de
macroeconomía. Son documentos recientes: conviene revisar la versión vigente
antes de citarlos.}

%------------------------------------------------------------------------------
\bloque{Bloque 6 --- Intermediación financiera y banca --- 3 semanas}
{Por qué existen los bancos, cómo gestionan sus riesgos y por qué son frágiles,
con la teoría de Freixas y Rochet. Cierra con el sistema financiero peruano.
Cubre la Unidad 4 de Sistemas Financieros y el primer tema de 1FIN31.}

\tema{6.1 Por qué existen los intermediarios financieros}
\lecturas{\ME{7} \quad Freixas y Rochet (1999), cap.~2}
\articulosB{Diamond (1984)}
\consultalib{\GTB{2--3} \quad Cecchetti y Schoenholtz (2015), cap.~11}
\consulta{Freixas y Rochet (1999), cap.~1}
\nota{Costos de transacción, información asimétrica, selección adversa y
riesgo moral; la «especialidad» de los bancos (\SA{1}). Diamond (1984)
explica al banco como monitor delegado.}

\tema{6.2 Bancos e instituciones financieras: balances y desempeño}
\lecturas{\SA{11--12}}
\consultalib{\GTB{4} \quad Saunders et al.\ (2022), caps.~14--19}
\nota{El balance de un banco comercial, el análisis de rentabilidad por el
ROE y una mirada a las demás instituciones: seguros, fondos de inversión, de
pensiones y empresas \emph{fintech}.}

\tema{6.3 Riesgos de las instituciones financieras}
\lecturas{\SA{20--23}}
\consultalib{\GTB{5--7} \quad Freixas y Rochet (1999), cap.~8}
\nota{Riesgo de crédito, de liquidez, de tasa de interés y de insolvencia, y su
gestión dentro del balance: brecha de duración, colchones de liquidez y
análisis de crédito.}

\tema{6.4 Depósitos, corridas, titulización y banca en la sombra}
\lecturas{\GTB{11--12} \quad Freixas y Rochet (1999), cap.~7}
\articulosB{Diamond y Dybvig (1983)}
\nota{El contrato de depósito, las corridas bancarias, el seguro de depósitos y
la intermediación fuera del sistema bancario. Diamond y Dybvig (1983) es el
modelo de referencia de las corridas; Freixas y Rochet lo extienden al riesgo
sistémico y al mercado interbancario.}

\tema{6.5 El sistema financiero peruano}
\lecturas{BCRP, \emph{Reporte de Estabilidad Financiera} (edición más reciente)}
\consultalib{Superintendencia de Banca, Seguros y AFP (SBS), estadísticas del sistema financiero}
\nota{Estructura del sistema: banca múltiple, financieras, cajas municipales y
rurales, AFP y seguros; concentración, dolarización del crédito y de los
depósitos, y morosidad. Quién supervisa qué (la SBS y el BCRP) y cómo funciona
el Fondo de Seguro de Depósitos. Conviene leerlo con los datos de la SBS a
mano: aplica al caso peruano los Bloques 1, 2 y 6.}

%------------------------------------------------------------------------------
\bloque{Bloque 7 --- Regulación, crisis y estabilidad financiera --- 3 semanas}
{Sigue el curso 1FIN31 (Regulación Financiera, P.~Collazos), que usa a Saunders
y a Freixas y Rochet, y el bloque de crisis de 1FIN27.}

\tema{7.1 Regulación bancaria: de Basilea I a Basilea III}
\lecturas{\SA{13} \quad \SA{13E} \quad \ME{18} \quad Freixas y Rochet (1999), cap.~9}
\articulos{Admati, DeMarzo, Hellwig y Pfleiderer (2011)}
\articulosB{Holmström y Tirole (1997)}
\consultalib{\GTB{15--16} \quad Hull (2015), caps.~15--16}
\consulta{Dewatripont y Tirole (1994)}
\nota{Por qué se regula a los bancos, el requerimiento de capital ponderado por
riesgo y su evolución, y la intervención del supervisor (modelo de
Dewatripont--Tirole). Holmström y Tirole (1997) muestran por qué el capital de
los bancos limita el crédito a las empresas.}

\tema{7.2 Crisis financieras: anatomía e historia}
\lecturas{\ME{8} \quad \GTB{14} \quad Kindleberger y Aliber, caps.~1--2}
\articulos{Kaminsky y Reinhart (1999) \quad Laeven y Valencia (2018)}
\consultalib{Kindleberger y Aliber, caps.~3--9 y 13}
\nota{El ciclo de expansión del crédito, burbuja y pánico; las crisis gemelas,
bancarias y de balanza de pagos; la crisis de 2008.}

\tema{7.3 Prestamista de última instancia y política macroprudencial}
\lecturas{Kindleberger y Aliber, caps.~11--12}
\articulos{Reinhart y Rogoff (2011) \quad Bernanke y Gertler (1995)}
\articulosB{Lorenzoni (2008) \quad Bianchi (2011)}
\consultalib{Jeanne y Korinek (2010) \quad Farhi y Werning (2016)}
\nota{El canal del crédito y el acelerador financiero conectan este tema con el
Bloque 15 del sílabo de macroeconomía (temas 15.4 y 15.5). Lorenzoni (2008) y
Bianchi (2011) explican por qué hay sobreendeudamiento: cada deudor no
internaliza el efecto de sus deudas sobre los precios de los activos que sirven
de colateral. Es la justificación de la política macroprudencial, y Bianchi la
desarrolla para economías emergentes con paradas súbitas.}

%------------------------------------------------------------------------------
\bloque{Bloque 8 --- Finanzas internacionales --- 2 semanas}
{Sigue el curso 1FIN27 (Finanzas Internacionales, A.~Segura): sistema
monetario internacional, crisis de economías emergentes y la política
monetaria en una economía semidolarizada.}

\tema{8.1 Tipo de cambio, paridad de tasas de interés y paridad del poder de compra}
\lecturas{\SA{9} \quad \ME{15} \quad \KO{14--16}}

\tema{8.2 El sistema monetario internacional y los regímenes cambiarios}
\lecturas{\ME{16} \quad \KO{18--19}}
\articulos{Obstfeld, Shambaugh y Taylor (2004) \quad Ilzetzki, Reinhart y Rogoff (2017)}
\nota{El patrón oro, Bretton Woods, la flotación y el trilema.}

\tema{8.3 Globalización financiera y crisis en economías emergentes}
\lecturas{\ME{25} \quad \KO{20} \quad \KO{22}}
\articulos{Calvo y Reinhart (2000) \quad Eichengreen, Hausmann y Panizza (2002)}
\consultalib{Uribe y Schmitt-Grohé (2015), caps.~6, 10 y 11}
\consulta{Eichengreen y Gupta (2016) \quad Rajan (2005)}
\nota{Miedo a flotar, pecado original, paradas súbitas y deuda soberana.}
\nota{Los capítulos de Uribe y Schmitt-Grohé siguen el borrador de marzo de 2015,
como en los sílabos de macroeconomía y de economía internacional: el 11 es la
deuda soberana, que en la edición publicada de 2017 es el capítulo 13.}

\tema{8.4 Política monetaria y cambiaria en una economía semidolarizada: el Perú}
\articulos{BCRP (2022) \quad Armas y Grippa (2006) \quad Rossini, Montoro y Luna (2020) \quad Morán (2017)}
\consulta{Montoro, Pérez y Herrada (2020)}
\consultalib{Rossini y Vega (2007) \quad Castillo, Pérez Forero y Tuesta (2011)}
\nota{Metas de inflación con dolarización parcial, encajes, intervención
cambiaria y \emph{swaps} del BCRP, y las medidas frente a la pandemia.}
\nota{Rossini y Vega (2007) y Castillo, Pérez Forero y Tuesta (2011) estudian el
mecanismo de transmisión en el Perú; los usa el curso de Política Monetaria del
BCRP. El tema 13.7 del sílabo de macroeconomía ve la misma pregunta desde la
teoría: metas de inflación con hojas de balance dolarizadas.}

%------------------------------------------------------------------------------
\newpage
\section*{Anexo A. De los capítulos a los bloques}
\addcontentsline{toc}{section}{Anexo A. De los capítulos a los bloques}

Capítulos de cada libro que entran en el plan y temas donde se leen. Lo que
queda fuera se detalla en el Anexo B.

\newcommand{\mapa}[2]{%
\subsection*{#1}
{\small\filasalternas
\noindent\begin{tabularx}{\textwidth}{>{\raggedright\arraybackslash}p{1.1cm}X>{\raggedright\arraybackslash}p{3.8cm}}
\toprule
\rowcolor{mhnavy}\cab{Cap.} & \cab{Título} & \cab{Temas} \\
\midrule
#2
\bottomrule
\end{tabularx}
}}

\mapa{A.1 Saunders, Cornett y Erhemjamts, \emph{Financial Markets and Institutions} (texto base)}{%
1 & Introduction & consulta (1.1) \\
2 & Determinants of Interest Rates & consulta (1.2) \\
3 & Interest Rates and Security Valuation (con apéndices 3A y 3B) & 1.3 \\
4 & The Federal Reserve System, Monetary Policy, and Interest Rates & consulta (1.4) \\
5--6 & Money Markets; Bond Markets & 2.1 \\
7 & Mortgage Markets & 2.3 \\
8 & Stock Markets (con apéndices 8A y 8B) & 2.2, 3.3 \\
9 & Foreign Exchange Markets & 8.1 \\
10 & Derivative Securities Markets & 2.4 \\
11--12 & Commercial Banks; Financial Statements and Analysis & 6.2 \\
13 & Regulation of Commercial Banks (con apéndice 13E) & 7.1 \\
20--23 & Types of Risks; Managing Credit, Liquidity, Interest Rate and Insolvency Risk & 6.3 \\
24 & Managing Risk off the Balance Sheet with Derivative Securities & 2.4 \\
25 & Loan Sales and Securitization & 2.3 \\
}

\mapa{A.2 Campbell, \emph{Financial Decisions and Markets}}{%
1 & Choice under Uncertainty & consulta (3.1) \\
2 & Static Portfolio Choice & 3.2 \\
3 & Static Equilibrium Asset Pricing & 3.3 \\
4 & The Stochastic Discount Factor (4.1--4.5) & 4.2, 4.3, 4.4 \\
5 & Present Value Relations & 5.1 \\
6 & Consumption-Based Asset Pricing & 5.2, 5.3 \\
7 & Production-Based Asset Pricing & 5.4 \\
8 & Fixed-Income Securities (8.1--8.4) & 5.5 \\
}

\mapa{A.3 Cochrane, \emph{Asset Pricing}}{%
1--2 & Consumption-Based Model and Overview; Applying the Basic Model & 4.1 \\
3--4 & Contingent Claims Markets; The Discount Factor & 4.2 \\
5--6 & Mean-Variance Frontier and Beta Representations; Relation between Discount Factors, Betas, and Mean-Variance Frontiers (selección) & 4.3 \\
9 & Factor Pricing Models (9.1--9.4) & 3.3 \\
10 & GMM in Explicit Discount Factor Models & 4.4 \\
12 & Regression-Based Tests of Linear Factor Models & 4.4 \\
19 & Term Structure of Interest Rates (19.1--19.3) & 5.5 \\
20 & Expected Returns in the Time Series and Cross Section (20.1) & 5.1 \\
21 & Equity Premium Puzzle and Consumption-Based Models & 5.2, 5.3 \\
}

\mapa{A.4 Mishkin y Eakins, \emph{Financial Markets and Institutions}}{%
1--2 & Why Study Financial Markets and Institutions?; Overview of the Financial System & 1.1 \\
3--5 & Interest Rates: Meaning, Changes, Risk and Term Structure & 1.2 (cap.~3, consulta en 1.3) \\
6 & Are Financial Markets Efficient? & 2.2 \\
7 & Why Do Financial Institutions Exist? & 6.1 \\
8 & Why Do Financial Crises Occur? & 7.2 \\
9--10 & Central Banks; Conduct of Monetary Policy & 1.4 \\
11--14 & Money, Bond, Stock and Mortgage Markets & consulta (2.1--2.3) \\
15--16 & The Foreign Exchange Market; The International Financial System & 8.1, 8.2 \\
18 & Financial Regulation & 7.1 \\
25 & Financial Crises in Emerging Market Economies & 8.3 \\
}

\mapa{A.5 Greenbaum, Thakor y Boot, \emph{Contemporary Financial Intermediation}}{%
2--3 & The Nature and Variety of Financial Intermediation; The What, How, and Why of Financial Intermediaries & consulta (6.1) \\
4 & Bank Risks & consulta (6.2) \\
5--7 & Interest Rate Risk; Liquidity Risk; Spot Lending and Credit Risk & consulta (6.3) \\
11--12 & Securitization; The Deposit Contract, Deposit Insurance, and Shadow Banking & 6.4 \\
14 & The 2007--2009 Financial Crisis and Other Financial Crises & 7.2 \\
15--16 & Objectives of Bank Regulation; Milestones in Banking Legislation and Regulatory Reform & consulta (7.1) \\
}

\mapa{A.6 Krugman, Obstfeld y Melitz, \emph{Economía internacional}}{%
14 & Los tipos de cambio y el mercado de divisas: un enfoque de activos & 8.1 \\
15 & El dinero, los tipos de interés y los tipos de cambio & 8.1 \\
16 & El nivel de precios y el tipo de cambio a largo plazo & 8.1 \\
18 & Los tipos de cambio fijos y la intervención en los mercados de divisas & 8.2 \\
19 & Sistemas monetarios internacionales: una visión histórica & 8.2 \\
20 & Globalización financiera: oportunidades y crisis & 8.3 \\
22 & Países en desarrollo: crecimiento, crisis y reforma & 8.3 \\
}

\mapa{A.7 Freixas y Rochet, \emph{Microeconomics of Banking}}{%
1 & Introduction & consulta (6.1) \\
2 & The Role of Financial Intermediaries & 6.1 \\
7 & Individual Bank Runs and Systemic Risk & 6.4 \\
8 & Managing Risks in the Banking Firm & consulta (6.3) \\
9 & The Regulation of Banks & 7.1 \\
}

%------------------------------------------------------------------------------
\section*{Anexo B. Lo que no entra en el plan}
\addcontentsline{toc}{section}{Anexo B. Lo que no entra en el plan}

{\small\filasalternas
\noindent\begin{tabularx}{\textwidth}{>{\raggedright\arraybackslash}p{4.6cm}X}
\toprule
\rowcolor{mhnavy}\cab{Materia} & \cab{Capítulos y secciones} \\
\midrule
Otras instituciones financieras & \SA{14--19} (consulta en el tema 6.2); \ME{19--22}, \ME{26--27} \\
Gestión de riesgos con fines profesionales & \ME{23--24}; \GTB{8--10}, \GTB{13}, \GTB{17--18} \\
Valoración de activos avanzada & \CA{9--12} (portafolio intertemporal, finanzas de los hogares, información asimétrica y liquidez); \CO{7--8}, \CO{11}, \CO{13--18}; \CO{20.2--20.3} \\
Macroeconomía abierta & \KO{13}, \KO{17}, \KO{21}; los capítulos de comercio internacional \\
Materias de otros cursos & Finanzas corporativas, gestión de portafolios, matemática financiera en tiempo continuo y valoración de derivados. Ver Anexo D \\
\bottomrule
\end{tabularx}
}

\medskip
\noindent La traducción de Freixas y Rochet de la carpeta (\emph{Economía
bancaria}) es una imagen de solo 29 páginas; el plan usa en su lugar la edición
en inglés de 1999 que está en la carpeta de macroeconomía, completa aunque sin
capa de texto. Las presentaciones de regulación financiera
(L01 a L13) también son imágenes sin texto (se revisaron L01 y L11).

%------------------------------------------------------------------------------
\section*{Anexo C. Programas utilizados}
\addcontentsline{toc}{section}{Anexo C. Programas utilizados}

{\small\filasalternas
\noindent\begin{tabularx}{\textwidth}{>{\raggedright\arraybackslash}p{6.0cm}X}
\toprule
\rowcolor{mhnavy}\cab{Programa} & \cab{Temas del plan} \\
\midrule
1ECO06, Introducción a las Finanzas (2025-2, A.~Segura) & Bloques 1 a 3 \\
1FIN32, Temas en Finanzas 1: mercado de capitales (2025-2, P.~Bringas) & Bloques 1 y 2; 3.2, 3.3 \\
1FIN09, Teoría y Gestión de Portafolios (2025-2, J.~Villavicencio) & Bloque 3; la gestión de portafolios, en el Anexo D.1 \\
1FIN31, Regulación Financiera (2025-2, P.~Collazos) & 6.1, 6.4, 7.1 \\
1FIN27, Finanzas Internacionales (2025-2, A.~Segura) & Bloques 7 y 8 \\
1F0124, Sistemas Financieros (Universidad del Pacífico, 2022-2) & Bloques 1, 2 y 6; 1.4 \\
\bottomrule
\end{tabularx}
}

\medskip
\noindent\destacado{Programas de la carpeta que no se usaron}, por pertenecer a
finanzas de empresa o a métodos ya cubiertos: Análisis Financiero 1 (1FIN02),
Finanzas Cuantitativas 1, 2 y Avanzadas (1FIN18, 1FIN04 y 1FIN25),
Instrumentos Financieros (1FIN06), Activos Alternativos (1FIN19), Finanzas
Corporativas 2 (FIN210), Microfinanzas (FIN220) y Plan de Negocios (GES235).
Tópicos de Econometría Financiera (FIN222) ya está en el material opcional del
sílabo de econometría. También se revisó el programa de Moneda y Banca de la
Universidad Nacional de Cajamarca, que sigue a Mishkin.

%------------------------------------------------------------------------------
\section*{Anexo D. Material opcional}
\addcontentsline{toc}{section}{Anexo D. Material opcional}

Todo lo de este anexo es opcional: no entra en el cómputo de horas ni modifica
el cronograma.

\subsection*{D.1 Finanzas de empresa y gestión de portafolios}
Berk y DeMarzo (2023); Bodie, Kane y Marcus (2018), partes V y VII; programas
1FIN09 y FIN210.

\subsection*{D.2 Valoración en tiempo continuo y derivados}
Sigue los programas de Finanzas Cuantitativas (1FIN18, 1FIN04 y 1FIN25): Hull
(2021); Back (2017); Duffie (2001); \CO{17--19}.

\subsection*{D.3 Medición de riesgos y econometría financiera}
Valor en riesgo, pérdida esperada y CoVaR: Hull (2015), caps.~12--14; Anexo D.1
del sílabo de macroeconometría. Modelos de volatilidad: Bloque 9 de ese
sílabo. Econometría de los mercados financieros: Campbell, Lo y MacKinlay
(1997).

\subsection*{D.4 Teoría de los mercados e instituciones desde los fundamentos}
Para profundizar el Bloque 6 con modelos: Sussman (2023); LeRoy y Werner
(2000); Hens y Rieger (2010), que añade finanzas conductuales.

\subsection*{D.5 \emph{Fintech}, criptomonedas y microfinanzas}
\SA{19}; el tema opcional de criptomonedas de 1FIN27; el programa FIN220
(Microfinanzas).

\subsection*{D.6 Proyecto aplicado con datos peruanos}
No tiene horas asignadas; conviene hacerlo en las sesiones de los sábados de
los Bloques 4 y 5.
\begin{enumerate}[leftmargin=1.6em,itemsep=2pt]
  \item \textbf{Curva de rendimiento.} Estimar la curva de los bonos soberanos en
        soles con el modelo de Nelson y Siegel (1987), con datos del BCRP, y
        leer su pendiente con la hipótesis de expectativas (tema 5.5).
  \item \textbf{CAPM y factores.} Con retornos de acciones de la Bolsa de
        Valores de Lima, estimar betas y aplicar regresiones de Fama y MacBeth
        (1973) (temas 3.3 y 4.4).
  \item \textbf{Ecuación de Euler.} Estimar por GMM la ecuación de Euler del
        consumo con datos peruanos, a la Hansen y Singleton (tema 4.4).
\end{enumerate}

%------------------------------------------------------------------------------
\section*{Anexo E. Referencias}
\addcontentsline{toc}{section}{Anexo E. Referencias}

\subsection*{E.1 Textos del plan}
{\small
\begin{itemize}[leftmargin=1.4em,itemsep=1pt,label={}]
\item Campbell, J. Y. (2017). \emph{Financial Decisions and Markets: A Course in Asset Pricing}. Princeton University Press.
\item Cochrane, J. H. (2005). \emph{Asset Pricing}, ed. revisada. Princeton University Press.
\item Freixas, X. y J.-C. Rochet (1999). \emph{Microeconomics of Banking}. MIT Press (carpeta de macroeconomía).
\item Greenbaum, S. I., A. V. Thakor y A. W. A. Boot (2019). \emph{Contemporary Financial Intermediation}, 4.ª ed. Academic Press.
\item Kindleberger, C. P. y R. Z. Aliber. \emph{Manías, pánicos y cracs: historia de las crisis financieras}. Versión de la carpeta, con epílogo 2010--2020.
\item Krugman, P. R., M. Obstfeld y M. J. Melitz (2016). \emph{Economía internacional: teoría y política}, 10.ª ed. Pearson.
\item Mishkin, F. S. y S. G. Eakins (2023). \emph{Financial Markets and Institutions}, 10.ª ed. global. Pearson.
\item Saunders, A., M. M. Cornett y O. Erhemjamts (2022). \emph{Financial Markets and Institutions}, 8.ª ed. McGraw-Hill.
\item BCRP. \emph{Reporte de Estabilidad Financiera}. Banco Central de Reserva del Perú, edición semestral más reciente.
\end{itemize}
}

\subsection*{E.2 Libros de consulta de la carpeta (sin horas)}
{\small
\begin{itemize}[leftmargin=1.4em,itemsep=1pt,label={}]
\item Back, K. E. (2017). \emph{Asset Pricing and Portfolio Choice Theory}, 2.ª ed. Oxford University Press.
\item Berk, J. y P. DeMarzo (2023). \emph{Corporate Finance}, 6.ª ed. Pearson.
\item Bodie, Z., A. Kane y A. J. Marcus (2018). \emph{Investments}, 11.ª ed. McGraw-Hill.
\item Campbell, J. Y., A. W. Lo y A. C. MacKinlay (1997). \emph{The Econometrics of Financial Markets}. Princeton University Press (carpeta de econometría).
\item Cecchetti, S. G. y K. L. Schoenholtz (2015). \emph{Money, Banking, and Financial Markets}, 4.ª ed. McGraw-Hill.
\item Duffie, D. (2001). \emph{Dynamic Asset Pricing Theory}, 3.ª ed. Princeton University Press.
\item Hens, T. y M. O. Rieger (2010). \emph{Financial Economics: A Concise Introduction to Classical and Behavioral Finance}. Springer.
\item Hull, J. C. (2015). \emph{Risk Management and Financial Institutions}, 4.ª ed. Wiley.
\item Hull, J. C. (2021). \emph{Options, Futures, and Other Derivatives}, 11.ª ed. Pearson.
\item LeRoy, S. F. y J. Werner (2000). \emph{Principles of Financial Economics}. Versión preliminar de los autores.
\item Sussman, O. (2023). \emph{The Economics of Financial Markets and Institutions: From First Principles}. Oxford University Press.
\item Tuckman, B. (2002). \emph{Fixed Income Securities}, 2.ª ed. Wiley.
\item Uribe, M. y S. Schmitt-Grohé (2015). \emph{Open Economy Macroeconomics}. Borrador del 17 de marzo de 2015; publicado por Princeton University Press, 2017 (carpeta de macroeconomía).
\end{itemize}
}

\subsection*{E.3 Referencias tomadas de los programas}
{\small
\begin{itemize}[leftmargin=1.4em,itemsep=1pt,label={}]
\item Admati, A., P. DeMarzo, M. Hellwig y P. Pfleiderer (2011). Fallacies, Irrelevant Facts, and Myths in the Discussion of Capital Regulation: Why Bank Equity is Not Expensive. Rock Center for Corporate Governance, Stanford University, Working Paper 86.
\item Armas, A. y F. Grippa (2006). Metas de inflación en una economía dolarizada: la experiencia del Perú. En A. Armas y A. Ize (eds.), \emph{Dolarización financiera: la agenda de política}. BCRP y FMI.
\item BCRP (2022). Estabilidad monetaria: diseño e implementación de la política monetaria. Folleto institucional.
\item Bernanke, B. y M. Gertler (1995). Inside the Black Box: The Credit Channel of Monetary Policy Transmission. NBER Working Paper 5146.
\item Calvo, G. y C. Reinhart (2000). Fear of Floating. NBER Working Paper 7993.
\item Dewatripont, M. y J. Tirole (1994). \emph{The Prudential Regulation of Banks}. MIT Press.
\item Eichengreen, B. y P. Gupta (2016). Managing Sudden Stops. World Bank Policy Research Working Paper 7639.
\item Eichengreen, B., R. Hausmann y U. Panizza (2002). Original Sin: The Pain, the Mystery and the Road to Redemption. Banco Interamericano de Desarrollo.
\item Ilzetzki, E., C. Reinhart y K. Rogoff (2017). Exchange Arrangements Entering the 21st Century: Which Anchor Will Hold? NBER Working Paper 23134.
\item Kaminsky, G. y C. Reinhart (1999). The Twin Crises: The Causes of Banking and Balance-of-Payments Problems. \emph{American Economic Review}, 89(3), 473--500.
\item Laeven, L. y F. Valencia (2018). Systemic Banking Crises Revisited. IMF Working Paper 18/206.
\item Montoro, C., F. Pérez y R. Herrada (2020). Medidas del BCRP frente a la pandemia del nuevo coronavirus. \emph{Revista Moneda}, 182, 10--18.
\item Morán, M. (2017). El swap cambiario del BCRP: ¿qué es y cómo funciona? \emph{Revista Moneda}, 170, 8--15.
\item Obstfeld, M., J. Shambaugh y A. Taylor (2004). The Trilemma in History: Tradeoffs among Exchange Rates, Monetary Policies and Capital Mobility. NBER Working Paper 10396.
\item Rajan, R. (2005). Has Financial Development Made the World Riskier? NBER Working Paper 11728.
\item Reinhart, C. y K. Rogoff (2011). From Financial Crash to Debt Crisis. \emph{American Economic Review}, 101(5), 1676--1706.
\item Rossini, R., C. Montoro y M. Luna (2020). Desarrollo del mercado financiero y política monetaria: la experiencia peruana. \emph{Revista Estudios Económicos}, BCRP, 39, 9--20.
\end{itemize}
}

\subsection*{E.4 Otras referencias}
{\small
\begin{itemize}[leftmargin=1.4em,itemsep=1pt,label={}]
\item Bianchi, J. (2011). Overborrowing and Systemic Externalities in the Business Cycle. \emph{American Economic Review}, 101(7), 3400--3426.
\item Castillo, P., F. Pérez Forero y V. Tuesta (2011). Los mecanismos de transmisión de la política monetaria en el Perú. \emph{Revista Estudios Económicos}, BCRP, 21.
\item Chow, T., B. Halperin y J. Z. Mazlish (2023). Transformative AI, Existential Risk, and Real Interest Rates. Documento de trabajo.
\item Diamond, D. W. (1984). Financial Intermediation and Delegated Monitoring. \emph{Review of Economic Studies}, 51(3), 393--414.
\item Diamond, D. W. y P. H. Dybvig (1983). Bank Runs, Deposit Insurance, and Liquidity. \emph{Journal of Political Economy}, 91(3), 401--419.
\item Fama, E. F. y J. D. MacBeth (1973). Risk, Return, and Equilibrium: Empirical Tests. \emph{Journal of Political Economy}, 81(3), 607--636.
\item Farhi, E. e I. Werning (2016). A Theory of Macroprudential Policies in the Presence of Nominal Rigidities. \emph{Econometrica}, 84(5), 1645--1704.
\item Hansen, L. P. y K. J. Singleton (1982). Generalized Instrumental Variables Estimation of Nonlinear Rational Expectations Models. \emph{Econometrica}, 50(5), 1269--1286.
\item Holmström, B. y J. Tirole (1997). Financial Intermediation, Loanable Funds, and the Real Sector. \emph{Quarterly Journal of Economics}, 112(3), 663--691.
\item Jeanne, O. y A. Korinek (2010). Excessive Volatility in Capital Flows: A Pigouvian Taxation Approach. \emph{American Economic Review}, 100(2), 403--407.
\item Lorenzoni, G. (2008). Inefficient Credit Booms. \emph{Review of Economic Studies}, 75(3), 809--833.
\item Lucas, R. E. (1978). Asset Prices in an Exchange Economy. \emph{Econometrica}, 46(6), 1429--1445.
\item Mehra, R. y E. C. Prescott (1985). The Equity Premium: A Puzzle. \emph{Journal of Monetary Economics}, 15(2), 145--161.
\item Nelson, C. R. y A. F. Siegel (1987). Parsimonious Modeling of Yield Curves. \emph{Journal of Business}, 60(4), 473--489.
\item Rossini, R. y M. Vega (2007). El mecanismo de transmisión de la política monetaria en un entorno de dolarización financiera: el caso del Perú entre 1996 y 2006. \emph{Revista Estudios Económicos}, BCRP, 14.
\item Van Nieuwerburgh, S. (2026). Financing the AI Buildout. \emph{Brookings Papers on Economic Activity}; NBER Working Paper 35865.
\end{itemize}
}

\vspace{1.2cm}
\franjas[2pt]
\begin{center}\small\color{mhslate}
Plan de lectura de elaboración propia. Los títulos de los capítulos
corresponden a las obras citadas y se reproducen únicamente con fines de
referencia bibliográfica. La asignación de capítulos y secciones se basa en
los índices de los libros de la carpeta; los temas de los cursos proceden de
los programas citados en el Anexo C.
\end{center}

\end{document}
